% in Anlehnung an https://github.com/derdanu/akad-vorlage
% Autor: TheFreakyL
%-------------------------------------------
% Vorgaben Assignment aus Studienheft SQL301
%-------------------------------------------
% Umfang: 8 - 10 Seiten (inkl. Abbildungen und Tabellen, aber ohne Deckblatt, % Gliederung und Literaturverzeichnis, Eidesstattliche Erklaerung)
% Zeilenabstand: 1,5, Tabellen 1,0 
% Schriftart: frei
% Schriftgrad: 11 oder 12 pt
% Tabellen auch 10pt
% Variablen, physikalische Groessen und Funktionszeichen werden kursiv gedruckt.
% Ränder: links: 4,5 cm, rechts 2,0 cm, oben und unten jeweils 3,0 cm
% Deckblatt: (Name, Adresse, AKAD-E-Mail-Adresse, Immatrikulationsnummer, Modulbezeichnung, Thema, Abgabedatum, Dozent)

% Reihenfolge (siehe SQL301, S. 71):
% Titelblatt, ohne Nummerierung
% Inhaltsverzeichnis (ohne Nummerierung) -> Abbildungsverzeichnis (Nummerierung in Römisch, beginnend mit "III")-> Tabellenverzeichnis -> Abk.-Verzeichniss
% Textteil, Nummerierung in arabisch, beginnen mit "1"
% Anhang -> Literaturverzeichnis, Nummerierung in Römisch, folgend auf Textteil

\documentclass[listof=totoc, bibliography=totoc, a4paper, 12pt]{scrartcl}
\usepackage[ngerman]{babel}
\usepackage{epsfig}
\usepackage{times}
\usepackage{tabto}
\usepackage{wrapfig}
\usepackage{multirow}
\usepackage[onehalfspacing]{setspace}
\usepackage{listings}
\usepackage{mathptmx}
\usepackage{geometry}
\usepackage{helvet}
\usepackage{courier}
\usepackage{setspace}
\usepackage{textcomp}
\usepackage[T1]{fontenc}
\usepackage[utf8]{inputenc}
\usepackage{fancyhdr}
\usepackage{float} % Notwendig fuer figure[h]
\usepackage[printonlyused]{acronym}
\usepackage{multicol} % für zweispaltiges Deckblatt

% Pakete für Literatur + zitieren
\usepackage[backend=biber, style=authoryear-ibid, autocite=footnote]{biblatex}
\usepackage{csquotes}
\addbibresource{Literatur/literatur.bib}

% Allgemeine Definitionen

% Assignmentbezeichnungen
%-----------------------------------------------------------------------------------
\newcommand*{\assi}{Assignment REG23-AS-002}
% Titel
\newcommand*{\titel}{Steuerungs- und Regelungstechnik} 
% Betreff
\newcommand*{\betreff}{Thema \glqq Bohreinrichtung\grqq{} } 
% Modul
\newcommand*{\modul}{REG23}
% Betreuer
\newcommand*{\dozent}{Heimerdinger} 
% Vor- und Nachname
\newcommand*{\name}{Papa Schlumpf}
% Straße und Hausnummer
\newcommand*{\strasse}{Am Pilz 01} 
% Plz und Ort
\newcommand*{\plzort}{1234 Schlumpfhausen} 
% Immatrikulationsnummer
\newcommand*{\immanr}{123456789}
% Studiengang
\newcommand{\stud}{Mechatronik - Bachelor of \tabto{40mm}Engineering (B. Eng.)}
%-----------------------------------------------------------------------------------

% URL in gleicher Schriftart wie Rest
\urlstyle{same}
% Email mit Verlinkung
\newcommand*{\email}{\href{mailto:hier_könnte_deine E-Mail_stehen@test.de}{hier_könnte_deine_Email_stehen@x.de}} 

% Formatierungen für das Literaturverzeichnis
%--------------------------------------------
% Doppelpunkt und Leerzeichen nach Autor
\renewcommand*{\labelnamepunct}{\addcolon\addspace} 
% Punkt am Ende entfernen
\renewcommand*{\finentrypunct}{\addspace}
% Komma zwischen einzelnen Einträgen
\renewcommand*{\newunitpunct}{\addcomma\addspace}
% Punkt am Ende von Zitaten entfernen
\renewcommand*{\bibfootnotewrapper}[1]{\bibsentence#1\addspace}
% Zeilenabstand zwischen einzelnen Einträgen
\setlength{\bibitemsep}{0.5\baselineskip}

% Paket für Verlinkungen im pdf-Dokument/ Einstellung für pdf-Dokument (z.B. Literaturverzeichnis, URLs, Email,...)
\usepackage[
%---------------------------------------------------------
% optional, Metadaten der .pdf
pdftitle={\titel},
pdfsubject={Assignment zum Modul REG23},
pdfauthor={\name},
pdfkeywords={akad, assignment, pdf, REG23, SPS, CODESYS}
%---------------------------------------------------------
% pdfborderstyle={/S/U/W 1}, % Links unterstrichen
pdfborder={0 0 0},  % Links nicht sichtbar im pdf
colorlinks = false,
pdfpagelabels,
pdfstartview = FitH,
bookmarksopen = true,
bookmarksnumbered = true,
linkcolor = black,
plainpages = false,
hypertexnames = false,
urlcolor = black,
citecolor = black] {hyperref}



\newcounter{savepage} % Zähler für Seitenzahl, für später ab Römischer Nummerierung
%--------------------------------------------------------------------------------------------------------------------------------
%                                       Dokument begin
%--------------------------------------------------------------------------------------------------------------------------------
\begin{document}

% Deckblatt
\newgeometry{left=20mm, right=20mm, top=30mm, bottom=30mm}
\thispagestyle{empty}
{\centering
\vspace*{\fill}
\Huge{\textbf{\assi}}

\vspace{\baselineskip}

\Huge{\titel}
\vspace{\baselineskip} \\
\Large{\betreff}

\vspace{\baselineskip}

\begin{multicols}{2}
\raggedright
\setlength{\columnseprule}{1pt}
\begin{small}

Name, Vorname: \tabto{40mm} \name \\
Immatrikulationsnr.:  \tabto{40mm} \immanr \\
Adressse: \tabto{40mm} \strasse \\
\tabto{40mm} \plzort \\
E-Mail: \tabto{40mm} \email

\columnbreak

Studiengang:  \tabto{40mm} \stud \\
Modul:  \tabto{40mm} \modul \\
Abgabe am:  \tabto{40mm} \today \\
Dozent:  \tabto{40mm} \dozent

\end{small}
\end{multicols}{}

\vfill\vfill

\includegraphics[scale=0.35]{Bilder/akad_logo.png}\par}



\newgeometry{left=45mm, right=20mm, top=30mm, bottom=30mm}

\begin{spacing}{1.0} % Verzeichnisse werden mit einzeiligem Abstand gesetzt

% Inhaltsverzeichnis ohne Nummerierung
\pagestyle{fancy}
\fancyhead{}
\renewcommand{\headrulewidth}{1pt}
\fancyhead{}
\fancyfoot{}

% Inhaltsverzeichnis
\tableofcontents 
\newpage

% ab hier Nummerierung oben in Römisch, beginnend mit 3
\pagestyle{fancy}
\fancyhead{}
\renewcommand{\headrulewidth}{1pt}
\fancyhead[R]{\thepage}
\fancyfoot{}
\pagenumbering{Roman}
\setcounter{page}{3}

% bei Nichtbedarf entsprechende Zeile Auskommentieren:

% Abbildungsverzeichnis, 
\listoffigures 
\newpage

% Tabellenverzeichnis
\listoftables 
\newpage

% Abkürzungsverzeichnis
\section*{Abkürzungsverzeichnis}
\addcontentsline{toc}{section}{Abkürzungsverzeichnis} 
\begin{acronym}
        \acro{SPS}{speicherprogrammierbare Steuerung}
\end{acronym}
\newpage

% Formelverzeichnis
%\listof{Formel}{Formelübersicht}
\end{spacing} 


% ab hier Nummerierung oben in arabisch, beginngend mit 1
\pagestyle{fancy}
\fancyhead{}
\renewcommand{\headrulewidth}{1pt}
\fancyhead[R]{\thepage}
\fancyfoot{}
\pagenumbering{arabic}

\begin{spacing}{1.5} % Zeilenabstand: 1,5 fuer den Textteil
%-------------------------------------------------------------------------------------------------------------------
%                     Ab hier der eigentliche Textteil
%-------------------------------------------------------------------------------------------------------------------
% Einleitung
\section{Einleitung}
TEST:\\
Im Zeitalter der Industrie 4.0 spielen Steuerungs- und Automatisierungssysteme eine zentrale Rolle. Daten der Anlagen werden erfasst und online ausgewertet. Erst diese Informationen machen es möglich, dass Maschinen autonom arbeiten können.\autocite[vgl.][24]{Seitz2021} Dieses Assignment gibt einen Einblick in die Grundlagen der Steuerungstechnik. Anschließend werden diese Grundlagen anhand einer \ac{SPS} dargestellt.
Das ist ein Test\autocite{oV2021a}


% Grundlagen
\section{Grundlagen der Steuerungstechnik}
\subsection{Sensoren}

\subsection{Aktoren}

\subsection{SPS}

\subsection{CODESYS}

\subsection{OSI-Schichtmodell}

% Hauppteil
\section{praktisches Beispiel}
\subsection{Lastenheft}

\subsection{Programmablaufplan}

\subsection{Programmierung}

\subsection{Simulation}

% Schluss
\section{Zusammenfassung}
Ziel der Arbeit ist es, ein Schrittkettensimulationsprogramm unter \ac{codesys} in \ac{fup} zu erstellen und zu simulieren. Hierzu sind im ersten Teil erforderliche Grundlagen dargestellt. Die in der Aufgabenstellung geforderte Einordnung der verwendeten Schichten des OSI-Modells bei einem Regelungssystem wird ebenso erläutert. Im Kapitel \textit{Bohreinrichtung} wird im ersten Schritt anhand der Aufgabenstellung ein Lastenheft erstellt. Darauf aufbauend wird der Programmablauf grafisch in einem Ablaufplan festgehalten. Dieser stellt die Grundlage für die eigentliche Programmierung dar. Zur Überprüfung der Funktionsfähigkeit des Codes wird anschließend eine Simulation der Schrittkette erstellt.\\
Die Grundlagen zu den Sensoren und Aktoren konnte aufgrund des Umfangs des Assignments nur sehr knapp beschrieben werden. Genauso wurde die Entwicklungsumgebung CoDeSys nur in seiner Grundform erläutert. Die gesamten Funktionen, die für die Programmierung erforderlich sind, konnten nicht beschrieben werden. An dieser Stelle ist auf das Benutzerhandbuch \autocite[][]{manCODESYS} zu verweisen. Die Simulation zeigt, dass die Programmierung den geforderten Schrittkettenablauf erfüllt. Das Ziel des Assign\-ments ist somit erreicht.\\
Aufgrund von fehlenden Spezifikationen entsprechend der Maschinenrichtlinien und nicht vorhandenen Sicherheitseinrichtungen für den Maschinenführer wird die Bohreinrichtung jedoch in der Industrie keine Anwendung finden. Dies sollte bei einer möglichen Weiterentwicklung umgesetzt werden. Zusätzlich ist eine Vorschub- und Drehzahlregelung empfehlenswert, um die Bohreinrichtung für verschiedene Werkstoffe und Bohrergrößen nutzen zu können.

\setcounter{savepage}{\value{page}} %speichert Seitenzahl für folgende Römische Nummerierung

\end{spacing}

\clearpage

% ab hier Römisch
\pagenumbering{Roman}
\setcounter{page}{\value{savepage}}

% Anhang, bei Nichtbedarf auskommentieren
%\newgeometry{left=30mm, right=10mm, top=30mm, bottom=20mm}
%Original: \newgeometry{left=45mm, right=20mm, top=30mm, bottom=30mm, landscape}

\section*{Anhang}
\addcontentsline{toc}{section}{Anhang}
\renewcommand{\thesubsection}{\Alph{subsection}}
%Begin Anhang
\subsection{Shared-Memory-API} \label{anhang:sharedMemApi}


\begin{landscape}
\subsection{Fehleruntersuchung}\label{anhang:error_analysis}
    %\input{Tabellen/error_analysis}
\end{landscape}

                    

%[frame=leftline,  fontsize=\scriptsize, breaklines=true, linenos, ]

    






   
%\clearpage

% Literaturverzeichnis
\printbibliography

\end{document}
